\documentclass[12pt]{article}
\usepackage[UTF8]{inputenc}
\usepackage{thaisupp}
\usepackage{enumitem}
\begin{document}
\title{แรงลอเรนตซ์}
\maketitle
\section*{In-class Questions}
\begin{enumerate}[label=\arabic*.]
\item แรงลอเรนตซ์คือแรงที่เกิดจากอะไร \hfill \textbf{\small (Main)}
\begin{enumerate}[label=)]
\item ก) ประจุไฟฟ้าที่อยู่นิ่งในสนามไฟฟ้า
\item ข) ประจุไฟฟ้าที่เคลื่อนที่ในสนามแม่เหล็ก
\item ค) ประจุไฟฟ้าที่เคลื่อนที่ในสนามไฟฟ้า
\item ง) มวลของอนุภาคที่เคลื่อนที่ในสนามโน้มถ่วง
\end{enumerate}
\item สูตรแรงลอเรนตซ์คือข้อใด \hfill \textbf{\small (Main)}
\begin{enumerate}[label=)]
\item ก) F = qvB
\item ข) F = qvB sinθ
\item ค) F = qvB tanθ
\item ง) F = qvB cosθ
\end{enumerate}
\item ทิศทางของแรงลอเรนตซ์หาได้โดยใช้กฎใด \hfill \textbf{\small (Main)}
\begin{enumerate}[label=)]
\item ก) กฎมือซ้าย
\item ข) กฎมือขวา
\item ค) กฎมือขวาของแอมแปร์
\item ง) กฎของฟาราเดย์
\end{enumerate}
\item เมื่อประจุบวกเคลื่อนที่เข้าสู่สนามแม่เหล็กในทิศตั้งฉากกับสนาม แรงลอเรนตซ์จะมีทิศทางอย่างไร \hfill \textbf{\small (Main)}
\begin{enumerate}[label=)]
\item ก) ขนานกับสนามแม่เหล็ก
\item ข) ตั้งฉากกับทั้งทิศทางความเร็วและสนามแม่เหล็ก
\item ค) ตรงข้ามกับทิศทางความเร็ว
\item ง) ขนานกับทิศทางความเร็ว
\end{enumerate}
\item ประจุลบที่เคลื่อนที่ในสนามแม่เหล็กจะได้รับแรงลอเรนตซ์ในทิศตรงข้ามกับประจุบวกหรือไม่ \hfill \textbf{\small (Main)}
\begin{enumerate}[label=)]
\item ก) ไม่ เพราะแรงขึ้นกับขนาดของประจุไม่ใช่ชนิด
\item ข) ใช่ เพราะทิศทางของแรงขึ้นกับชนิดของประจุ
\item ค) ไม่ เพราะแรงลอเรนตซ์เป็นปริมาณสเกลาร์
\item ง) ขึ้นกับทิศทางของสนามแม่เหล็กเท่านั้น
\end{enumerate}
\item เมื่อประจุเคลื่อนที่เข้าสู่สนามแม่เหล็กในทิศขนานกับสนาม ประจุจะเคลื่อนที่อย่างไร \hfill \textbf{\small (Main)}
\begin{enumerate}[label=)]
\item ก) เป็นวงกลม
\item ข) เป็นเกลียว
\item ค) เป็นเส้นตรง
\item ง) หยุดนิ่ง
\end{enumerate}
\item เมื่อประจุเคลื่อนที่ตั้งฉากกับสนามแม่เหล็ก ประจุจะเคลื่อนที่เป็นรูปร่างใด \hfill \textbf{\small (Main)}
\begin{enumerate}[label=)]
\item ก) เส้นตรง
\item ข) พาราโบลา
\item ค) วงกลม
\item ง) วงรี
\end{enumerate}
\item รัศมีของวงโคจรการเคลื่อนที่เป็นวงกลมของประจุในสนามแม่เหล็ก ขึ้นอยู่กับปริมาณใด \hfill \textbf{\small (Main)}
\begin{enumerate}[label=)]
\item ก) มวลและความเร็วของประจุเท่านั้น
\item ข) ประจุและความเร็วของประจุเท่านั้น
\item ค) มวล ประจุ ความเร็ว และสนามแม่เหล็ก
\item ง) มวลและสนามแม่เหล็กเท่านั้น
\end{enumerate}
\item อิเล็กตรอนมีประจุ 1.6 × 10⁻¹⁹ C เคลื่อนที่ด้วยความเร็ว 10⁶ m/s ตั้งฉากกับสนามแม่เหล็ก 0.1 T แรงลอเรนตซ์ที่กระทำต่ออิเล็กตรอนมีค่าเท่าใด \hfill \textbf{\small (Main)}
\begin{enumerate}[label=)]
\item ก) 1.6 × 10⁻¹⁴ N
\item ข) 1.6 × 10⁻¹³ N
\item ค) 1.6 × 10⁻¹² N
\item ง) 1.6 × 10⁻¹⁵ N
\end{enumerate}
\item โปรตอนเคลื่อนที่เป็นวงกลมในสนามแม่เหล็ก โดยมีรัศมีวงโคจร 0.2 เมตร มีความเร็ว 10⁵ m/s ถ้าสนามแม่เหล็กเพิ่มเป็น 2 เท่า รัศมีวงโคจรจะเปลี่ยนเป็นเท่าใด \hfill \textbf{\small (Main)}
\begin{enumerate}[label=)]
\item ก) 0.05 เมตร
\item ข) 0.1 เมตร
\item ค) 0.4 เมตร
\item ง) 0.8 เมตร
\end{enumerate}
\item มวลสเปกโทรมิเตอร์ (mass spectrometer) ใช้หลักการใดในการแยกไอโซโทป \hfill \textbf{\small (Main)}
\begin{enumerate}[label=)]
\item ก) ความแตกต่างของประจุไฟฟ้า
\item ข) ความแตกต่างของมวล
\item ค) ความแตกต่างของสนามไฟฟ้า
\item ง) ความแตกต่างของสภาพร้อน
\end{enumerate}
\item ในมวลสเปกโทรมิเตอร์ ไอโซโทปที่มีมวลมากกว่าจะตกบนฟิล์มที่ระยะทางจากจุดเข้ามากกว่าหรือน้อยกว่า \hfill \textbf{\small (Main)}
\begin{enumerate}[label=)]
\item ก) น้อยกว่า
\item ข) มากกว่า
\item ค) เท่ากัน
\item ง) ขึ้นกับประจุของไอโซโทป
\end{enumerate}
\item ไซโคลตรอน (cyclotron) ทำงานบนหลักการใด \hfill \textbf{\small (Main)}
\begin{enumerate}[label=)]
\item ก) การเคลื่อนที่เป็นวงกลมของประจุในสนามแม่เหล็กและการเพิ่มพลังงานโดยสนามไฟฟ้า
\item ข) การสั่นพ้องของคลื่นแม่เหล็กไฟฟ้า
\item ค) การชนของอนุภาคในสนามไฟฟ้า
\item ง) การหมุนของแม่เหล็กรอบแกน
\end{enumerate}
\item ข้อจำกัดของไซโคลตรอนคืออะไร \hfill \textbf{\small (Main)}
\begin{enumerate}[label=)]
\item ก) ไม่สามารถเร่งอิเล็กตรอนได้
\item ข) ไม่สามารถทำงานได้ถ้าความเร็วใกล้ความเร็วแสง
\item ค) ไม่สามารถเร่งประจุบวกได้
\item ง) ไม่สามารถทำงานในสุญญากาศ
\end{enumerate}
\item อนุภาคอัลฟา (He²⁺) เคลื่อนที่ด้วยความเร็ว 10⁵ m/s ตั้งฉากกับสนามแม่เหล็ก 0.2 T อนุภาคอัลฟาจะเคลื่อนที่เป็นวงกลมรัศมีเท่าใด (มวลอนุภาคอัลฟา = 6.64 × 10⁻²⁷ kg) \hfill \textbf{\small (Main)}
\begin{enumerate}[label=)]
\item ก) 1.04 × 10⁻² m
\item ข) 1.04 × 10⁻³ m
\item ค) 1.04 × 10⁻¹ m
\item ง) 1.04 × 10⁻⁴ m
\end{enumerate}
\item เมื่อประจุเคลื่อนที่ในสนามแม่เหล็กในมุม θ ≠ 0°, 90°, 180° ประจุจะเคลื่อนที่เป็นรูปร่างใด \hfill \textbf{\small (Main)}
\begin{enumerate}[label=)]
\item ก) วงกลม
\item ข) พาราโบลา
\item ค) เกลียว
\item ง) เส้นตรง
\end{enumerate}
\item โปรตอนเคลื่อนที่ด้วยความเร็ว 2 × 10⁶ m/s ตั้งฉากกับสนามแม่เหล็ก 0.5 T จงหาคาบการโคจรของโปรตอน (มวลโปรตอน = 1.67 × 10⁻²⁷ kg, q = 1.6 × 10⁻¹⁹ C) \hfill \textbf{\small (Main)}
\begin{enumerate}[label=)]
\item ก) 1.31 × 10⁻⁷ s
\item ข) 1.31 × 10⁻⁶ s
\item ค) 1.31 × 10⁻⁸ s
\item ง) 1.31 × 10⁻⁵ s
\end{enumerate}
\item ในมวลสเปกโทรมิเตอร์ ถ้าไอโซโทปของธาตุ X มีประจุ +1 และมวล 28 u อีกไอโซโทปหนึ่งมีประจุ +1 และมวล 30 u เข้าสู่สนามแม่เหล็กด้วยความเร็วเท่ากัน รัศมีวงโคจรของไอโซโทปมวล 28 u เป็นกี่เท่าของไอโซโทปมวล 30 u \hfill \textbf{\small (Main)}
\begin{enumerate}[label=)]
\item ก) 28/30 เท่า
\item ข) 30/28 เท่า
\item ค) 30²/28² เท่า
\item ง) √(28/30) เท่า
\end{enumerate}
\item ในไซโคลตรอน ความถี่ของสนามไฟฟ้าที่ใช้เร่งประจุต้องเท่ากับข้อใด \hfill \textbf{\small (Main)}
\begin{enumerate}[label=)]
\item ก) ความถี่เชิงมุมของการโคจร
\item ข) ความถี่เชิงมุมของการโคจรคูณ 2
\item ค) ความถี่เชิงมุมของการโคจรหาร 2
\item ง) ความถี่เชิงมุมของการโคจรหาร π
\end{enumerate}
\item ถ้าความเร็วของประจุเพิ่มขึ้นเป็น 2 เท่า แรงลอเรนตซ์จะเปลี่ยนแปลงอย่างไร \hfill \textbf{\small (Main)}
\begin{enumerate}[label=)]
\item ก) เพิ่มเป็น 2 เท่า
\item ข) เพิ่มเป็น 4 เท่า
\item ค) เพิ่มเป็น √2 เท่า
\item ง) ลดลงเป็น 1/2 เท่า
\end{enumerate}
\item อิเล็กตรอนเคลื่อนที่เข้าสู่สนามแม่เหล็ก 0.1 T ด้วยความเร็ว 10⁷ m/s จงหารัศมีวงโคจร (มวลอิเล็กตรอน = 9.11 × 10⁻³¹ kg, q = 1.6 × 10⁻¹⁹ C) \hfill \textbf{\small (Main)}
\begin{enumerate}[label=)]
\item ก) 5.69 × 10⁻⁴ m
\item ข) 5.69 × 10⁻³ m
\item ค) 5.69 × 10⁻⁵ m
\item ง) 5.69 × 10⁻² m
\end{enumerate}
\item เมื่อใช้กฎมือซ้ายของเฟลมมิง นิ้วหัวแม่มือ นิ้วชี้ และนิ้วกลาง แสดงถึงปริมาณใดตามลำดับ \hfill \textbf{\small (Main)}
\begin{enumerate}[label=)]
\item ก) แรง ความเร็ว สนามแม่เหล็ก
\item ข) สนามแม่เหล็ก แรง ความเร็ว
\item ค) แรง สนามแม่เหล็ก ความเร็ว
\item ง) ความเร็ว แรง สนามแม่เหล็ก
\end{enumerate}
\item ไซโคลตรอนสามารถเร่งอนุภาคได้ถึงพลังงานสูงสุดเท่าใดโดยประมาณ \hfill \textbf{\small (Main)}
\begin{enumerate}[label=)]
\item ก) หลายร้อย MeV
\item ข) หลายพัน MeV
\item ค) หลายสิบ MeV
\item ง) หลายร้อย GeV
\end{enumerate}
\item ความถี่เชิงมุมของการโคจรของประจุในสนามแม่เหล็ก (cyclotron frequency) มีค่าเท่าใด \hfill \textbf{\small (Main)}
\begin{enumerate}[label=)]
\item ก) ω = qB/m
\item ข) ω = mB/q
\item ค) ω = qm/B
\item ง) ω = B/qm
\end{enumerate}
\item สายไฟตรงยาว 0.5 เมตร มีกระแส 10 A ไหลในทิศตั้งฉากกับสนามแม่เหล็ก 0.2 T แรงลอเรนตซ์ที่กระทำต่อสายไฟมีค่าเท่าใด \hfill \textbf{\small (Main)}
\begin{enumerate}[label=)]
\item ก) 1 N
\item ข) 2 N
\item ค) 0.5 N
\item ง) 4 N
\end{enumerate}
\item แรงลอเรนตซ์มีค่าสูงสุดเมื่อใด \hfill \textbf{\small (Extra)}
\begin{enumerate}[label=)]
\item ก) เมื่อประจุเคลื่อนที่ขนานกับสนามแม่เหล็ก
\item ข) เมื่อประจุเคลื่อนที่ตั้งฉากกับสนามแม่เหล็ก
\item ค) เมื่อประจุหยุดนิ่ง
\item ง) เมื่อประจุเคลื่อนที่ในทิศตรงข้ามกับสนามแม่เหล็ก
\end{enumerate}
\item แรงลอเรนตซ์ทำงาน (work) หรือไม่ บนประจุที่เคลื่อนที่ในสนามแม่เหล็ก \hfill \textbf{\small (Extra)}
\begin{enumerate}[label=)]
\item ก) ทำงานเสมอ
\item ข) ไม่ทำงานเพราะแรงตั้งฉากกับการกระจัด
\item ค) ทำงานเฉพาะเมื่อประจุเคลื่อนที่ขนานกับสนาม
\item ง) ทำงานเฉพาะเมื่อประจุเคลื่อนที่ตั้งฉากกับสนาม
\end{enumerate}
\item ประจุบวกเคลื่อนที่เข้าสู่สนามแม่เหล็กที่ชี้ลงในแนวดิ่ง ถ้าประจุเคลื่อนที่ไปทางทิศเหนือ แรงลอเรนตซ์จะมีทิศทางใด \hfill \textbf{\small (Extra)}
\begin{enumerate}[label=)]
\item ก) ทิศตะวันออก
\item ข) ทิศตะวันตก
\item ค) ทิศใต้
\item ง) ทิศบน
\end{enumerate}
\item ถ้าสนามแม่เหล็กเพิ่มขึ้น 3 เท่า แรงลอเรนตซ์จะเปลี่ยนแปลงอย่างไร \hfill \textbf{\small (Extra)}
\begin{enumerate}[label=)]
\item ก) เพิ่มขึ้น 3 เท่า
\item ข) เพิ่มขึ้น 9 เท่า
\item ค) ลดลง 3 เท่า
\item ง) เพิ่มขึ้น √3 เท่า
\end{enumerate}
\item เมื่อประจุเคลื่อนที่เป็นวงกลมในสนามแม่เหล็ก แรงลอเรนตซ์ทำหน้าที่เป็นอะไร \hfill \textbf{\small (Extra)}
\begin{enumerate}[label=)]
\item ก) แรงสู่ศูนย์กลาง
\item ข) แรงเสียดทาน
\item ค) แรงต้าน
\item ง) แรงย้อนกลับ
\end{enumerate}
\item การเปลี่ยนแปลงของปริมาณใดไม่มีผลต่อรัศมีวงโคจรของประจุในสนามแม่เหล็ก \hfill \textbf{\small (Extra)}
\begin{enumerate}[label=)]
\item ก) มวลของประจุ
\item ข) ความเร็วของประจุ
\item ค) ประจุของอนุภาค
\item ง) สนามแม่เหล็ก
\end{enumerate}
\item ประจุ 2 ชนิดมีอัตราส่วน q/m เท่ากัน เคลื่อนที่ด้วยความเร็วเท่ากันในสนามแม่เหล็กเดียวกัน ข้อใดถูกต้อง \hfill \textbf{\small (Extra)}
\begin{enumerate}[label=)]
\item ก) มีรัศมีวงโคจรเท่ากัน
\item ข) มีคาบการโคจรเท่ากัน
\item ค) มีแรงลอเรนตซ์เท่ากัน
\item ง) มีทิศทางการเคลื่อนที่เดียวกัน
\end{enumerate}
\item อิเล็กตรอนและโปรตอนเคลื่อนที่ด้วยความเร็วเท่ากันในสนามแม่เหล็กเดียวกัน ข้อใดถูกต้องเกี่ยวกับรัศมีวงโคจร \hfill \textbf{\small (Extra)}
\begin{enumerate}[label=)]
\item ก) รัศมีของอิเล็กตรอนใหญ่กว่า
\item ข) รัศมีของโปรตอนใหญ่กว่า
\item ค) รัศมีเท่ากัน
\item ง) ขึ้นกับทิศทางการเคลื่อนที่
\end{enumerate}
\item ในมวลสเปกโทรมิเตอร์แบบเบเบต (Bainbridge mass spectrometer) อนุภาคถูกเร่งโดยสนามไฟฟ้าก่อนเข้าสู่สนามแม่เหล็ก เพื่อให้อนุภาคมีอะไรเท่ากัน \hfill \textbf{\small (Extra)}
\begin{enumerate}[label=)]
\item ก) ความเร็ว
\item ข) พลังงานจลน์
\item ค) โมเมนตัม
\item ง) คาบการโคจร
\end{enumerate}
\item สายไฟยาว 0.3 เมตร มีกระแส 5 A ไหลในทิศตั้งฉากกับสนามแม่เหล็ก 0.4 T แรงลอเรนตซ์ที่กระทำต่อสายไฟมีค่าเท่าใด \hfill \textbf{\small (Extra)}
\begin{enumerate}[label=)]
\item ก) 0.6 N
\item ข) 6 N
\item ค) 60 N
\item ง) 0.06 N
\end{enumerate}
\item ข้อใดไม่ใช่ข้อจำกัดของไซโคลตรอน \hfill \textbf{\small (Extra)}
\begin{enumerate}[label=)]
\item ก) ผลสัมพัทธภาพทำให้มวลเพิ่มขึ้น
\item ข) ความถี่เรโซแนนซ์เปลี่ยน
\item ค) สามารถเร่งอิเล็กตรอนได้ยาก
\item ง) ต้องใช้สนามแม่เหล็กคงที่
\end{enumerate}
\item โปรตอนเคลื่อนที่ในสนามแม่เหล็ก 0.5 T ด้วยความเร็ว 10⁶ m/s จงหารัศมีวงโคจร (m = 1.67 × 10⁻²⁷ kg, q = 1.6 × 10⁻¹⁹ C) \hfill \textbf{\small (Extra)}
\begin{enumerate}[label=)]
\item ก) 2.09 × 10⁻² m
\item ข) 2.09 × 10⁻³ m
\item ค) 2.09 × 10⁻¹ m
\item ง) 2.09 × 10⁻⁴ m
\end{enumerate}
\item อนุภาคอัลฟา (q = 3.2 × 10⁻¹⁹ C, m = 6.64 × 10⁻²⁷ kg) เคลื่อนที่ด้วยความเร็ว 2 × 10⁶ m/s ตั้งฉากกับสนามแม่เหล็ก 0.1 T แรงลอเรนตซ์ที่กระทำมีค่าเท่าใด \hfill \textbf{\small (Extra)}
\begin{enumerate}[label=)]
\item ก) 6.4 × 10⁻¹⁴ N
\item ข) 6.4 × 10⁻¹³ N
\item ค) 6.4 × 10⁻¹² N
\item ง) 6.4 × 10⁻¹⁵ N
\end{enumerate}
\item ความถี่เชิงมุมของการโคจร ω = qB/m มีหน่วยเป็นอะไร \hfill \textbf{\small (Extra)}
\begin{enumerate}[label=)]
\item ก) m/s
\item ข) rad/s
\item ค) Hz
\item ง) 1/s²
\end{enumerate}
\item ในมวลสเปกโทรมิเตอร์ ถ้าต้องการแยกไอโซโทป ²³⁵U⁺ และ ²³⁸U⁺ ที่มีประจุ +1 เท่ากัน ไอโซโทปใดจะตกบนฟิล์มใกล้จุดเข้ามากกว่า \hfill \textbf{\small (Extra)}
\begin{enumerate}[label=)]
\item ก) ²³⁵U⁺
\item ข) ²³⁸U⁺
\item ค) ตกที่จุดเดียวกัน
\item ง) ขึ้นกับความเร็ว
\end{enumerate}
\item สายไฟพับสองเส้นยาว 0.2 เมตร/เส้น วางในสนามแม่เหล็ก 0.5 T ทำมุม 30° กับสนาม มีกระแส 3 A ไหล จงหาแรงลอเรนตซ์รวม \hfill \textbf{\small (Extra)}
\begin{enumerate}[label=)]
\item ก) 0.3 N
\item ข) 0.15 N
\item ค) 0.6 N
\item ง) 0.45 N
\end{enumerate}
\item อิเล็กตรอนเคลื่อนที่เป็นวงกลมในสนามแม่เหล็ก ถ้าสนามแม่เหล็กเพิ่มขึ้น 4 เท่า คาบการโคจรจะเปลี่ยนแปลงอย่างไร \hfill \textbf{\small (Extra)}
\begin{enumerate}[label=)]
\item ก) เพิ่มเป็น 4 เท่า
\item ข) ลดลงเป็น 1/4 เท่า
\item ค) เพิ่มเป็น 2 เท่า
\item ง) ลดลงเป็น 1/2 เท่า
\end{enumerate}
\item ถ้าความถี่ของไซโคลตรอนคือ 10 MHz สำหรับโปรตอน ความถี่สำหรับดิวเทอรอนในสนามแม่เหล็กเดียวกันจะเป็นเท่าใด (ดิวเทอรอนมี q = e, m ≈ 2mₚ) \hfill \textbf{\small (Extra)}
\begin{enumerate}[label=)]
\item ก) 20 MHz
\item ข) 10 MHz
\item ค) 5 MHz
\item ง) 40 MHz
\end{enumerate}
\item ถ้าโปรตอนและอิเล็กตรอนเคลื่อนที่ด้วยพลังงานจลน์เท่ากัน เข้าสู่สนามแม่เหล็กตั้งฉาก ข้อใดถูกต้อง \hfill \textbf{\small (Extra)}
\begin{enumerate}[label=)]
\item ก) รัศมีวงโคจรเท่ากัน
\item ข) รัศมีวงโคจรของอิเล็กตรอนเล็กกว่า
\item ค) รัศมีวงโคจรของโปรตอนเล็กกว่า
\item ง) คาบการโคจรเท่ากัน
\end{enumerate}
\item อนุภาคมีประจุ +2e เคลื่อนที่ในสนามแม่เหล็ก B = 0.2 T ด้วยความเร็ว 10⁵ m/s ทำมุม 30° กับสนาม จงหารัศมีของการเคลื่อนที่เป็นเกลียว \hfill \textbf{\small (Extra)}
\begin{enumerate}[label=)]
\item ก) 5.19 × 10⁻⁴ m
\item ข) 5.19 × 10⁻³ m
\item ค) 5.19 × 10⁻⁵ m
\item ง) 5.19 × 10⁻² m
\end{enumerate}
\item สายไฟวางในสนามแม่เหล็กโดยกระแสไฟฟ้าในสายไฟทำมุม 45° กับสนามแม่เหล็ก ถ้าสายไฟยาว 1 เมตร มีกระแส 10 A และสนามแม่เหล็ก 0.2 T แรงลอเรนตซ์มีค่าเท่าใด \hfill \textbf{\small (Extra)}
\begin{enumerate}[label=)]
\item ก) √2 N
\item ข) 1 N
\item ค) 2 N
\item ง) 2√2 N
\end{enumerate}
\item โปรตอนเคลื่อนที่ในสนามแม่เหล็กเดียวกันสองครั้ง ครั้งแรกมีความเร็ว v ครั้งที่สองมีความเร็ว 2v อัตราส่วนรัศมีวงโคจร r₂/r₁ เป็นเท่าใด \hfill \textbf{\small (Extra)}
\begin{enumerate}[label=)]
\item ก) 1
\item ข) 2
\item ค) 1/2
\item ง) 4
\end{enumerate}
\item ไซโคลตรอนใช้เร่งอนุภาคชนิดใดได้ดีที่สุด \hfill \textbf{\small (Extra)}
\begin{enumerate}[label=)]
\item ก) อิเล็กตรอน
\item ข) โปรตอนและไอออนหนัก
\item ค) โฟตอน
\item ง) นิวตริโน
\end{enumerate}
\item ถ้าประจุเคลื่อนที่ในสนามแม่เหล็กและสนามไฟฟ้าพร้อมกัน แรงลอเรนตซ์รวมคือข้อใด \hfill \textbf{\small (Extra)}
\begin{enumerate}[label=)]
\item ก) F = qE + qvB
\item ข) F = qE × qvB
\item ค) F = qE + qvB sinθ
\item ง) F = q(E + vB)
\end{enumerate}
\item อิเล็กตรอนเคลื่อนที่ในสนามแม่เหล็ก 0.01 T จงหาความถี่การโคจร (f = qB/(2πm)) \hfill \textbf{\small (Extra)}
\begin{enumerate}[label=)]
\item ก) 2.8 × 10⁸ Hz
\item ข) 2.8 × 10⁶ Hz
\item ค) 2.8 × 10⁷ Hz
\item ง) 2.8 × 10⁵ Hz
\end{enumerate}
\item สายไฟ 2 เส้นวางขนานกันห่างกัน 0.1 เมตร มีกระแส 5 A ในทิศเดียวกันในสนามแม่เหล็กภายนอก 0.1 T ที่ตั้งฉากกับสายไฟทั้งสอง แรงลอเรนตซ์ต่อความยาว 1 เมตรของสายไฟแต่ละเส้นมีค่าเท่าใด \hfill \textbf{\small (Extra)}
\begin{enumerate}[label=)]
\item ก) 0.5 N/m
\item ข) 0.25 N/m
\item ค) 0.75 N/m
\item ง) 1.0 N/m
\end{enumerate}
\item การเคลื่อนที่เป็นเกลียวของประจุในสนามแม่เหล็ก ระยะห่างระหว่างแต่ละรอบ (pitch) ของเกลียวขึ้นอยู่กับอะไร \hfill \textbf{\small (Extra)}
\begin{enumerate}[label=)]
\item ก) ความเร็วในทิศขนานกับสนามและคาบการโคจร
\item ข) ความเร็วในทิศตั้งฉากกับสนามเท่านั้น
\item ค) มวลของประจุเท่านั้น
\item ง) ประจุของอนุภาคเท่านั้น
\end{enumerate}
\item ในมวลสเปกโทรมิเตอร์ ถ้าใช้ความต่างศักย์เร่ง 1000 V สำหรับประจุ +1 อนุภาคที่มีประจุ +2 จะต้องใช้ความต่างศักย์เท่าใดเพื่อให้ได้ความเร็วเท่ากัน \hfill \textbf{\small (Extra)}
\begin{enumerate}[label=)]
\item ก) 250 V
\item ข) 500 V
\item ค) 2000 V
\item ง) 4000 V
\end{enumerate}
\item เมื่อใช้กฎมือซ้ายของเฟลมมิงกับสายไฟที่มีกระแสไหล ถ้าสนามแม่เหล็กชี้ลง และกระแสไหลไปทางทิศเหนือ แรงลอเรนตซ์จะชี้ไปทางใด \hfill \textbf{\small (Extra)}
\begin{enumerate}[label=)]
\item ก) ทิศตะวันออก
\item ข) ทิศตะวันตก
\item ค) ทิศขึ้น
\item ง) ทิศลง
\end{enumerate}
\item โปรตอนเคลื่อนที่ด้วยพลังงานจลน์ 1 MeV (1 MeV = 1.6 × 10⁻¹³ J) ในสนามแม่เหล็ก 0.5 T ตั้งฉากกับสนาม จงหารัศมีวงโคจร (m = 1.67 × 10⁻²⁷ kg) \hfill \textbf{\small (Extra)}
\begin{enumerate}[label=)]
\item ก) 0.129 m
\item ข) 0.258 m
\item ค) 0.0645 m
\item ง) 0.0129 m
\end{enumerate}
\end{enumerate}
\end{document}
